\section{Conclusions}
\label{sec:conclusions}

This study presents \dataset, a large-sample hydrological dataset for \ntotal Russian catchments following the CAMELS framework, covering \years.
The key contributions are:

\textbf{1.~Dataset composition and quality control.}
\dataset consolidates \ndischarge discharge gauges and \nwaterlevel water level gauges with systematic quality control.
The high-quality subset of \nhighquality gauges achieves median data completeness of 97.5\%.
Manual expert verification of every watershed boundary (\ntotal total) achieved mean areal errors of \meanerror against official Roshydromet data.

\textbf{2.~Hydrological signatures.}
Analysis of \nsigngauges discharge records yields \nsignatures hydrological signatures spanning wide ranges: mean discharge \meandischarge (0.017--8.346~mm/day), baseflow index \meanbaseflowindex (0.210--0.903), and mean half-flow date in early May (day~\meanhalfflowday).

\textbf{3.~Meteorological forcing consistency.}
Three precipitation products (ERA5-Land: \erafiveannual; MSWEP: \mswepannual; GPCP: \gpcpannual) show pairwise daily correlations ranging from $r = \eragpcpcorr$ to $\eramsweprcorr$, with systematic differences in annual totals.
ERA5-Land yields $Q/P \leq 1$ (median \medianrunoffratio) across all catchments, though this partly reflects ERA5-Land's high-latitude wet bias rather than independent validation (see Sect.~\ref{sec:technical_validation}).

Key limitations are discussed in Sect.~\ref{sec:limitations}.
To our knowledge, \dataset is the first publicly available CAMELS-standard dataset for the Russian Federation.
The dataset enables benchmarking of hydrological models across the largest previously unrepresented landmass in the global CAMELS network.

The dataset is publicly available under \license license at \placeholder{[DOI: 10.XXXX/zenodo.XXXXXXX]}.
